\scanpage{89}
% Typescript inverted in the scan.
\uline{Remarque} 3.13. Voici une généralisation parfois utile de 3.5. b). Considérons
une partie $\mathrm{E}$ de $\mathrm{X}$ telle que pour tout morphisme étale $\mathrm{f:X'\quad X}$, $\mathrm{E'=f^{-1}(E)}$
soit ``très dense'' dans $\mathrm{X'}$ (EGA IV 10.1.3.)\deleted{\unsure{x}}, i.e. que pour tout $\mathrm{f}$ comme ci-dessus,
tout ouvert $\mathrm{U}$ de $\mathrm{X'}$ qui contient $\mathrm{E'}$ soit égal à $\mathrm{X'}$. Alors les foncteurs
fibres $\mathrm{F_{\bar x}}$ sur $\mathrm{X_{\text{ét}}}$, pour $\mathrm{x\quad E}$, forment déjà une ``famille conservative''
(et par suite les corollaires 3.6. , 3.7. et 3.8. restent valables
en se bornant à y choisir $\mathrm{x\quad E}$). Cela se voit immédiatement en reprenant la
démonstration donnée de 3.5. b). Le résultat précédent s'applique notamment
dans les siutations suivantes:

a) $\mathrm{X}$ est un préschéma de Ja\deleted{\unsure{v|x}}cobson (EGA IV 10.4.1.), par exemple localement
de type fini sur un corps, ou sur $\mathrm{\operatorname{Spec}(\underline{Z})}$, et $\mathrm{E}$ est l'ensemble des
points fermés de $\mathrm{X}$.

b) $\mathrm{X}$ est localement de type fini sur un préschéma $\mathrm{S}$, et $\mathrm{E}$ est l'ensemble
des points de $\mathrm{X}$ qui sont fermé dans leur fibre.

c) $\mathrm{X}$ est localement noethérien, et $\mathrm{E}$ est l'ensemble des $\mathrm{x\quad X}$ tels que
$\mathrm{\bar x}$ soit un ensemble fini (i.e. un schéma semi-local de dimension $\mathrm{\quad 1}$) cf
EGA IV 10.5.3. et \deleted{\unsure{1|l}}10.5.5.
% archive sheet number: 88

